\articleappendix{Gamma 函数若干值的乘积的代数性（附录）}{Algebraicity of some products of values of the $\Gamma$ function}{N. Koblitz 与 A. Ogus}{2}{343--346}

令
\[
 A_N=N^{-1}\Z/\Z-\{0\},
\]
令 $U_N=(\Z/N\Z)^*$ 通过乘法作用于 $A_N$。对函数 $f:A_N\longrightarrow\C$，定义 $\langle f\rangle:U_N\longrightarrow\C$ 为
\[
 \langle f\rangle(u)=\sum_{a\in A_N}\langle a\rangle f(ua),
\]
其中 $\langle a\rangle$ 表示 $a$ 在 $0$ 与 $1$ 之间的代表元（representative）。

我们先计算
\[
 H_{\Q}=\{f:A_N\longrightarrow\Q:\langle f\rangle\text{ 为常数}\}.
\]

\paragraph*{例（Example）.}
若 $n$ 是 $N$ 的因子，$a\in A_N$，考虑集合
\[
 S_{n,a}=\bigl[\{a+k/n:k=0,1,\ldots,n-1\}\cup\{-na\}\bigr]\cap A_N.
\]
容易验证，它的示性函数（characteristic function）$\varepsilon_{n,a}$ 属于 $H_{\Q}$。\footnote{\textbf{译者注（本次校订）：}这里必须保留重数：应取 $\varepsilon_{n,a}=\sum_{k=0}^{n-1}\delta_{a+k/n}+\delta_{-na}$，再限制到 $A_N$，其中 $\delta_b$ 是点 $b$ 的示性函数。若取普通集合的示性函数，断言不成立。例如 $N=6,n=2,a=1/3$ 时，$S_{n,a}=\{1/3,5/6\}$，相应的 $\langle\varepsilon_{n,a}\rangle$ 在 $u=1,5$ 处分别为 $7/6,5/6$。按重数计数后两值均为 $3/2$，也与下文的 $(p+1)/2$ 一致。同理，当 $a=1/2$ 时应有 $\varepsilon_{1,a}=2\delta_a$。}注意，$S_{1,a}=\{a,-a\}$。

\begin{originaltheorem}{命题（Proposition）}
$H_{\Q}$ 由下列函数生成：
\[
 \{\varepsilon_{n,a}:n=1\text{，或 }n\text{ 为素数且 }na\ne0\}.
\]
\end{originaltheorem}
\begin{proof}
$U_N$ 在 $A_N$ 上作用的轨道（orbit）与满足 $1<d\leq N$ 的 $N$ 的因子 $d$ 对应：每个轨道唯一写成 $U_N(1/d)$ 的形式。$1/d$ 的稳定子群（stabilizer subgroup）为
\[
 I_d=\{u\in U_N:u\equiv1\pmod d\},
\]
而 $1/d$ 的轨道与 $U_d$ 典范同构（canonically isomorphic）。因此，$A_N$ 上的函数 $f$ 由一族函数 $f_d:U_d\longrightarrow\Q$ 决定，其中
\[
 f_d(v)=f(v/d).
\]

证明命题时，先将空间复化（complexification），以便使用 $U_N$ 的特征（character）。

\begin{originaltheorem}{引理（Lemma）}
若 $f:A_N\longrightarrow\C$，而 $\chi$ 是 $U_N$ 的一个特征，则 $\langle f\rangle$ 与 $\chi$ 的内积（inner product）为
\[
 (\langle f\rangle\mid\chi)
 =-\sum_d L(0,\chi_d)|I_d|\cdot(f_d\mid\chi_d),
\]
其中求和遍历满足以下条件的因子 $d\mid N$：$d>1$，且 $\chi$ 是 $U_d$ 上某个特征 $\chi_d$ 的拉回（pull-back）。\footnote{\textbf{译者注（本次校订）：}本引理须加上 $\chi\ne1$ 的条件，下文也只在这一情形使用它。对平凡特征，$d>1$ 时 $L(0,1_d)=0$，而左端实际等于 $|U_N|\sum_{a\in A_N}f(a)/2$，一般不为零。原因是公式 $L(0,\chi_d)=-d^{-1}\sum_{v\in U_d}|v|\chi_d(v)$ 仅适用于非平凡特征；平凡特征还带有常数项 $\varphi(d)/2$。}
\end{originaltheorem}

我们有
\begin{align*}
 (\langle f\rangle\mid\chi)
 &=\sum_{u\in U_N}\langle f\rangle(u)\overline{\chi}(u)\\
 &=\sum_{a\in A_N}\sum_{u\in U_N}\langle a\rangle f(ua)\overline{\chi}(u)\\
 &=\sum_{d\mid N}\sum_{v\in U_d}\sum_{u\in U_N}
   \langle v/d\rangle f_d(u_dv)\overline{\chi}(u),
\end{align*}
其中 $u_d\in U_d$ 为 $u$ 的像。现在选取 $U_N\longrightarrow U_d$ 的一组陪集代表元（coset representatives）$U'_d\subseteq U_N$，则
\begin{align*}
 \sum_{u\in U_N}f_d(u_dv)\overline\chi(u)
 &=\sum_{u'\in U'_d}\sum_{u\in I_d}
       f_d(u'_dv)\overline\chi(u')\overline\chi(u)\\
 &=\sum_{u'\in U'_d}f_d(u'_dv)\overline\chi(u')
       \sum_{u\in I_d}\overline\chi(u).
\end{align*}
若 $\chi$ 在 $I_d$ 上非平凡，该和便为零；若 $\chi$ 是 $U_d$ 上某个特征 $\chi_d$ 的拉回，该和则等于
\begin{align*}
 \sum_{u'\in U_d}f_d(u'v)\overline\chi_d(u')|I_d|
 &=\sum_{w\in U_d}f_d(w)\overline\chi_d(w)\chi_d(v)|I_d|\\
 &=\chi_d(v)|I_d|(f_d\mid\chi_d).
\end{align*}
将其代入上面的 $(\langle f\rangle\mid\chi)$ 表达式，并注意
$\langle v/d\rangle=|v|/d$，其中 $|v|$ 是 $v\in U_d$ 的最小正代表元，可得
\begin{align*}
 (\langle f\rangle\mid\chi)
 &=\sum_d\sum_{v\in U_d}|v|\chi_d(v)|I_d|d^{-1}(f_d\mid\chi_d)\\
 &=-\sum_d L(0,\chi_d)|I_d|(f_d\mid\chi_d),
\end{align*}
这正是所要证明的引理。

回到命题的证明。令 $d_f$ 为满足 $f_{d_f}\ne0$ 的最大因子 $d_f\mid N$。只需证明，每个 $f\in H$ 均能写成 $g+f'$，其中 $g$ 为这些 $\varepsilon$ 的线性组合（linear combination），并且 $d_{f'}<d_f$。令 $d=d_f$。我们将在下面证明，$f_d$ 是下述两类函数的线性组合：一类为经由
\[
 U_d\longrightarrow U_d/\{\pm1\}
\]
分解的函数 $h_1$；另一类为经由
\[
 U_d\longrightarrow U_{d/p}
\]
分解的函数 $h_p$，其中 $p$ 是 $d$ 的某个素因子且 $p\ne d$。第一类函数在 $\pm1$ 的作用下不变，因此在轨道 $U_N(1/d)$ 上相应的函数是一些 $\varepsilon_{1,a}$ 的线性组合。第二类函数在 $U_d\longrightarrow U_{d/p}$ 的核（kernel）$K$ 的作用下不变；若 $p\mid d/p$，则该核的阶为 $p$，否则为 $p-1$。第一种情形下，
\[
 K=\{1+kd/p:k=0,\ldots,p-1\};
\]
第二种情形下，则从这个集合中删去一个元素，即能被 $p$ 整除的那个 $1+kd/p$。因此，集合 $S_{p,w/d}$ 是 $(Kw)/d$ 再加上轨道 $U_N(p/d)$ 中的一个或两个元素。于是，轨道 $U_N(1/d)$ 上相应的函数显然可以写成一些 $\varepsilon_{p,a}$ 与支集（support）包含在 $p/d$ 的轨道中的函数的线性组合，从而得到了所需分解 $f=g+f'$。

最后证明上述关于 $f_d$ 的断言。若 $d$ 是某个奇数的两倍，则映射 $U_d\longrightarrow U_{d/2}$ 是同构，断言显然成立。在其余情形，选取 $U_d$ 的一个本原奇特征（primitive odd character）$\chi_d$，并令 $\chi$ 为 $\chi_d$ 在 $U_N$ 上的拉回。由于 $\chi$ 非平凡且 $\langle f\rangle$ 为常数，故
$(\langle f\rangle\mid\chi)=0$。由引理给出的内积公式，非零项只能来自满足 $f_e\ne0$ 且 $I_e\subseteq\ker(\chi)$ 的那些 $f_e$。但此时 $\ker(\chi)$ 同时包含 $I_e$ 和 $I_d$，因而也包含 $I_dI_e=I_m$，其中 $m=\gcd(d,e)$。由于 $\chi$ 是本原的，$m=d$，从而 $d\mid e$；又因 $e>d$ 时 $f_e=0$，所以
\[
 0=(\langle f\rangle\mid\chi)
   =-L(0,\chi_d)|I_d|(f_d\mid\chi_d).
\]
$\chi_d$ 是奇的且本原，故推出 $(f_d\mid\chi_d)=0$，即 $f_d$ 与 $U_d$ 的每个本原奇特征正交（orthogonal）。由此可知，$f_d$ 是偶特征（even character，即经由 $U_d/\{\pm1\}$ 分解）与非本原特征（imprimitive character，即经由某个 $U_{d/p}$ 分解）的线性组合。当 $d=p$ 时，还应注意，$U_p$ 上的常值函数也在 $\pm1$ 下不变，因此已包含于第一种情形。
\end{proof}

\begin{originaltheorem}{定理（Theorem）}
设 $f:A_N\longrightarrow\Q$ 满足 $\langle f\rangle$ 为常数，即对于所有 $u\in U_N$，有
\[
 \sum_a\langle a\rangle f(ua)=k.
\]
则
\[
 \Gamma(f)\mathrel{\mathop:}=\pi^{-k}\prod_a\Gamma(\langle a\rangle)^{f(a)}
\]
是代数数（algebraic number）。
\end{originaltheorem}
\begin{proof}
将任意 $f$ 映到 $\sum_a\langle a\rangle f(a)$ 的映射
$\langle\ \rangle:H_{\Q}\longrightarrow\Q$ 显然是线性的，因而对 $H_{\Q}$ 中的 $f,g$，有
\[
 \Gamma(f+g)=\Gamma(f)\Gamma(g).
\]
由命题，$H_{\Q}$ 中的任意 $f$ 都是这些 $\varepsilon$ 的 $\Q$ 线性组合；清去分母后，存在正整数 $n$，使 $nf$ 成为这些 $\varepsilon$ 的整数线性组合。\footnote{\textbf{译者注（本次校订）：}这里的整数系数可以为负，因此应使用“整数线性组合”。由于 $\Gamma(f)$ 非零，负系数只对应代数数的倒数，后续论证仍然成立。}由于 $\Gamma(f)^n=\Gamma(nf)$，只需对命题中描述的那些 $\varepsilon$ 验证定理。注意，$\langle\varepsilon_{1,a}\rangle=1$，且当 $pa\ne0$ 时，$\langle\varepsilon_{p,a}\rangle=(p+1)/2$；于是只须应用下述恒等式：\footnote{译注：原书下面第二个恒等式印作求和符号，且下限为 $k=1$；这一形式与通常的 Gamma 乘法公式不同，疑为排印错误。通常的公式为 $\prod_{k=0}^{d-1}\Gamma(x+k/d)=d^{1/2-dx}(2\pi)^{(d-1)/2}\Gamma(dx)$。下式依原书保留，以便区分原文与译注。}
\[
 \Gamma(x)\Gamma(1-x)=\pi(\sin\pi x)^{-1},
\]
\[
 \sum_{k=1}^{p-1}\Gamma\left(x+\frac{k}{p}\right)
 =p^{1/2-x}(2\pi)^{(p-1)/2}\Gamma(px).
\]
\end{proof}

\paragraph*{注（Remark）.}
Kubert 最近用另一套术语给出了一个更精确的结果。由其结果可知，若 $H_{\Q}$ 中的函数 $f$ 取整数值，则 $2f$ 是这些 $\varepsilon$ 的 $\Z$ 线性组合（即将发表）。

% 原书345页Gamma乘法公式显示为求和号且下限k=1；依扫描保留，未静默改动。

\origsection{}{Deligne 正文及附录的共同参考文献（Bibliographie）}
\begin{thebibliography}{SGA $4\frac12$}
\bibitem[1]{c36-1} M. V. Borovoi, \textit{Sur l'action du groupe de Galois sur les classes de cohomologie rationnelles de type $(p,p)$ des variétés abéliennes}, Mat. Sb. \textbf{94} (1974), 649--652（俄文）.
\bibitem[2]{c36-2} J. Coates and S. Lichtenbaum, \textit{On $l$-adic zeta functions}, Ann. of Math. (2) \textbf{98} (1973), 498--550.
\bibitem[3]{c36-3} P. Deligne, \textit{Applications de la formule des traces aux sommes trigonométriques}, SGA $4\frac12$, 168--232.
\bibitem[4]{c36-4} P. Deligne, \textit{Les constantes des équations fonctionnelles des fonctions $L$}, Modular Functions of One Variable. II, Lecture Notes in Math., vol. 349, Springer-Verlag, New York, 1973, 501--595.
\bibitem[5]{c36-5} P. Deligne, \textit{Formes modulaires et représentations $l$-adiques}, Séminaire Bourbaki 355（1969年2月）, Lecture Notes in Math., vol. 179, Springer-Verlag, New York, 139--172.
\bibitem[6]{c36-6} P. Deligne, \textit{Théorie de Hodge. III}, Inst. Hautes Études Sci. Publ. Math. \textbf{44} (1974), 5--77.
\bibitem[7]{c36-7} B. Gross, \textit{On the periods of abelian integrals and a formula of Chowla and Selberg}, Invent. Math. \textbf{45} (1978), 193--211.
\bibitem[8]{c36-8} Ju. I. Manin, \textit{Correspondences, motifs and monoidal transformations}.
\bibitem[9]{c36-9} Ju. I. Manin, \textit{Points paraboliques et fonction zêta des courbes modulaires}, Izv. \textbf{36} (1972), 19--66（俄文）.
\bibitem[10]{c36-10} I. I. Piateckii-Shapiro, \textit{Relations entre les conjectures de Tate et de Hodge pour les variétés abéliennes}, Mat. Sb. \textbf{85}(4) (1971), 610--620; Math. USSR-Sb. \textbf{14} (1971), 615--625.
\bibitem[11]{c36-11} J. P. Serre, \textit{Abelian $l$-adic representations and elliptic curves} (McGill), Benjamin, New York, 1968.
\bibitem[12]{c36-12} J. P. Serre, \textit{Facteurs locaux des fonctions zêta des variétés algébriques (définitions et conjectures)}, Sém. Delange--Pisot--Poitou 1969/70, exposé 19.
\bibitem[13]{c36-13} G. Shimura, \textit{On some arithmetic properties of modular forms of one and several variables}, Ann. of Math. (2) \textbf{102} (1975), 491--515.
\bibitem[14]{c36-14} C. L. Siegel, \textit{Berechnung von Zetafunktionen an ganzzahligen Stellen}, Nachr. Akad. Wiss. Göttingen Math. Phys. Kl. II \textbf{10} (1969), 87--102.
\bibitem[15]{c36-15} J. Tate, \textit{On the conjecture of Birch and Swinnerton-Dyer and a geometric analogue}, Séminaire Bourbaki 306 (1965/66), Benjamin, New York; 重印于 \textit{10 exposés sur la théorie des schémas}, North-Holland, Amsterdam, 1968.
\bibitem[16]{c36-16} A. Weil, \textit{Jacobi sums as grossencharaktere}, Trans. Amer. Math. Soc. \textbf{73} (1952), 487--495.
\bibitem[17]{c36-17} A. Weil, \textit{Sommes de Jacobi et caractères de Hecke}, Nachr. Akad. Wiss. Göttingen Math. Phys. Kl. II \textbf{1} (1974), 1--14.
\bibitem[18]{c36-18} D. Zagier, \textit{Modular forms whose Fourier coefficients involve zeta functions of quadratic fields}, Modular Forms of One Variable. VI, Lecture Notes in Math., vol. 627, Springer-Verlag, New York, 1977, 105--169.
\bibitem[19]{c36-19} N. Saavedra, \textit{Catégories tannakiennes}, Lecture Notes in Math., vol. 265, Springer-Verlag, New York, 1972.
\bibitem[SGA]{c36-sga} \textit{Séminaire de géométrie algébrique du Bois-Marie}.
\bibitem[SGA 4]{c36-sga4} M. Artin, A. Grothendieck et J. L. Verdier, \textit{Théorie des topos et cohomologie étale des schémas}, Lecture Notes in Math., vols. 259, 270, 305, Springer-Verlag, New York.
\bibitem[SGA $4\frac12$]{c36-sga45} P. Deligne, \textit{Cohomologie étale}, Lecture Notes in Math., vol. 569, Springer-Verlag, New York, 1977.
\end{thebibliography}
\begin{flushright}Institut des Hautes Études Scientifiques, Bures-sur-Yvette, France\end{flushright}
